\section{Discussion}
\paragraph{When does the 37\% rule actually apply?} Our sweeps make the
conditions concrete. The $1/e$ cutoff is close to optimal only when four things
hold at once: the number of options is fixed and reasonably large; you observe
ranks but care only about getting the single best; looking is free; and there is
no recall. Loosen any one and the threshold moves --- toward $0$ under search
cost, and toward the destination (hold out longer) as good options become
common in the parking problem. Treating $37\%$ as a universal heuristic for
apartments, dates, or parking spots is therefore usually wrong; the quantity
that transfers is the \emph{method} (compute the continuation value and stop
when the current option beats it), not the number.

\paragraph{What the agents tell us.} The value-based agent recovers the optimal
policy because its tabular representation can express the true, state-dependent
threshold, and because exact scoring lets us read off regret rather than a
noisy return. The policy-gradient agent's difficulty is instructive: a single
learned cutoff is the right hypothesis class for the stationary secretary rule
but cannot express the parking optimum at all. ``Learning the rule from
experience'' thus decomposes into two requirements --- enough data to estimate
the continuation values, and a policy class rich enough to hold the answer.

\paragraph{Limitations.} The problems are deliberately small so that exact
ground truth exists; we do not claim these particular learners scale to
high-dimensional stopping tasks. The parking model assumes independent
occupancy and a known street; correlated availability and unknown length are
natural next steps. The residual $Q$-learning regret in the sparse-occupancy
regime reflects finite-sample estimation of the acceptance boundary, not a
failure of the method, and shrinks under the larger-budget configuration.
